\Needspace{18\baselineskip}
\section{Zero-exponent limits on products with zero lower mass}
\label{sec:nonlp-selection}

Let $h$ be a Hermitian form on $X$. For a density $f$ with positive threshold, \cref{cor:individual-positive-threshold-exponential} shows that the normalization scalars $c_\lambda$ of the greatest exponential solutions $u_\lambda$ converge to $b_h(f)$ and that their Monge--Amp\`ere measures converge in total variation. Two questions about the potentials remain: whether the normalized solutions $v_\lambda=u_\lambda-\sup_Xu_\lambda$ converge as $\lambda\downarrow0$, and whether the shifted solutions
\[
 u_\lambda-\lambda^{-1}\log b_h(f)=v_\lambda+\lambda^{-1}\log\bigl(c_\lambda/b_h(f)\bigr)
\]
converge. The second question concerns the rescaled scalar error $\lambda^{-1}\log(c_\lambda/b_h(f))$, which also measures the loss in the lower barrier $v_\lambda\geq v-\lambda^{-1}\log(c_\lambda/b_h(f))$ for normalized threshold subsolutions $v$ (\cref{sec:individual-exponential-equations}). For $L^p$ densities with $p>1$, both questions have positive answers by \cref{thm:vanishing-exponential-selection}. We answer them for the integrable densities of \cref{thm:zero-lower-mass-product-equation}, which belong to no $L^p$ with $p>1$ and live on a background of zero lower mass.

The greatest solutions for these densities are pulled back from the projective factor and are invariant under the unitary group fixing the pole. For an invariant potential, the Monge--Amp\`ere measure of a ball around the pole determines the radial derivative of the potential. Convergence of the measures in total variation therefore becomes uniform convergence of radial derivatives, and the whole family of normalized solutions converges. The shifted solutions converge exactly when the limiting potential is integrable against the density. The limit is then the translate of the threshold solution whose mean against the density vanishes; otherwise the rescaled scalar error tends to $+\infty$. Both conclusions persist when the density varies together with the exponent.

We keep the notation of \cref{sec:zero-lower-mass-product-equations} with $Y=\mathbb P^k$, $k\geq1$: $(X_0,h_0)$ is a connected compact Hermitian manifold of dimension $m$ with $v_-(h_0)=0$ and $V_0=\int_{X_0}h_0^m$; $X=X_0\times Y$ has dimension $N=m+k$ and background $H=h_0+\kappa$; $C=\binom Nm$, $\mu=V_0^{-1}h_0^m\wedge\kappa^k$ and $R=\mu/\mathrm{d}V$; and $o$, $\psi$, $u_\alpha$ and $q_\alpha$ are as in \cref{lem:projective-sublogarithmic-density}. Let $\pi_Y\colon X\to Y$ be the projection, and put
\[
 b=CV_0,\qquad f_\alpha=Rq_\alpha\qquad(0<\alpha<1),
\]
so that $f_\alpha\,\mathrm{d}V=q_\alpha\mu$ and $b_H(f_\alpha)=b$ by \cref{thm:zero-lower-mass-product-equation}. As in \cref{subsec:zero-threshold-product-exponential}, put $L=1-\psi\geq1$ and $t=\mathrm{e}^{1-L}$ on $Y\setminus\{o\}$, and for $0<\alpha<1$ let
\[
 Q_\alpha=t+(1-t)\alpha L^{\alpha-1},\qquad D_\alpha=-\frac{\mathrm{d}Q_\alpha}{\mathrm{d}L}
\]
be the functions $Q$ and $D$ of that subsection. By \cref{lem:projective-exponential-barrier-growth} and its proof, for $x\geq1$,
\begin{equation}\label{eq:nonlp-selection-radial-data}
 \int_{\{L\geq x\}}q_\alpha\kappa^k=Q_\alpha(x)^k,\qquad
 \int_{\{L\geq x\}}\kappa^k=\mathrm{e}^{-k(x-1)},
\end{equation}
and $q_\alpha=Q_\alpha^{k-1}D_\alpha\mathrm{e}^{k(L-1)}$ on $Y\setminus\{o\}$. For $T\geq1$ the compact sets $Y_T=\{L\leq T\}$ exhaust $Y\setminus\{o\}$, and $Y_1$ is the hyperplane $\{z_0=0\}$. On the affine chart $\{z_0\neq0\}$, with coordinates $z\in\C^k$, we put $s=\log|z|^2$; then $L=1-s+\log(1+\mathrm{e}^s)$ is a decreasing function of $s$. The \emph{first moment} of the threshold solution $u_\alpha$ is
\begin{equation}\label{eq:nonlp-selection-moment}
 J_\alpha=\int_Y(-u_\alpha)q_\alpha\kappa^k
         =\int_X(-\pi_Y^*u_\alpha)f_\alpha\,\mathrm{d}V\in(0,+\infty].
\end{equation}

The unitary group $G=\mathrm U(k)$ acts on $Y$ by $g\cdot[z_0:z']=[z_0:gz']$, where $z'=(z_1,\ldots,z_k)$. It fixes $o$ and preserves the hyperplane $\{z_0=0\}$, the form $\kappa$ and the function $\psi$. Hence it preserves $u_\alpha$, and also $q_\alpha$, since $q_\alpha\kappa^k=\mathrm{MA}^{\mathrm{ext}}_\kappa(u_\alpha)$. A function on $Y$ is called \emph{invariant} if it is invariant under $G$.

\subsection{Invariant potentials on projective space}

For an invariant potential, the computation of \cref{lem:projective-sublogarithmic-density} expresses the Monge--Amp\`ere measure of the balls centered at $o$ through the radial derivative of the potential. We need it for potentials that are neither smooth nor bounded.

\Needspace{14\baselineskip}
\begin{lemma}\label{lem:nonlp-selection-radial}\index{Radial equation!invariant potentials on projective space}
Let $v\in\PSH(Y,\kappa)$ be invariant. There is a convex nondecreasing function $\Phi_v\colon\mathbb R\to\mathbb R$, whose right derivative $\Phi_v'$ satisfies $0\leq\Phi_v'\leq1$, such that $v=\Phi_v(s)-\log(1+\mathrm{e}^s)$ on $\C^k\setminus\{0\}$. Write $\Phi_v'(\pm\infty)$ for the limits of $\Phi_v'$ at $\pm\infty$. Then $\mathrm{MA}^{\mathrm{ext}}_\kappa(v)$ charges neither $o$ nor the hyperplane $\{z_0=0\}$, and
\begin{equation}\label{eq:nonlp-selection-flux}
 \mathrm{MA}^{\mathrm{ext}}_\kappa(v)\bigl(\{0<|z|^2\leq\mathrm{e}^s\}\bigr)
 =\Phi_v'(s)^k-\Phi_v'(-\infty)^k\qquad(s\in\mathbb R).
\end{equation}
In particular, its total mass is $\Phi_v'(+\infty)^k-\Phi_v'(-\infty)^k\leq1$. Suppose moreover that $v\geq u_\delta-c$ for some $0<\delta<1$ and some constant $c$. Then $\mathrm{MA}^{\mathrm{ext}}_\kappa(v)$ is a probability measure,
\[
 \Phi_v'(s)=\mathrm{MA}^{\mathrm{ext}}_\kappa(v)\bigl(\{|z|^2\leq\mathrm{e}^s\}\bigr)^{1/k}\qquad(s\in\mathbb R),
\]
and $v$ is continuous on $Y\setminus\{o\}$. Finally, two invariant functions in $\PSH(Y,\kappa)$ whose extended Monge--Amp\`ere measures are the same probability measure differ by a constant.
\end{lemma}
\begin{proof}
On the affine chart, $\kappa=\ddc\log(1+|z|^2)$, so $p=v+\log(1+|z|^2)$ is plurisubharmonic on $\C^k$. It is invariant under $\mathrm U(k)$, which acts transitively on the spheres centered at the origin; hence $p=\Phi_v(\log|z|^2)$ off the origin for a function $\Phi_v$ on $\mathbb R$. Restricted to a complex line through the origin, $p$ is subharmonic, and it is not identically $-\infty$ there, since otherwise invariance would make $p$ identically $-\infty$ off the origin. By \cref{cor:radial-subharmonic-convexity}, $\Phi_v$ is finite and convex. It is bounded above at $-\infty$, because $p$ is bounded above near the origin; a convex function with this property is nondecreasing. Since $v\leq\sup_Yv$, we have $\Phi_v(s)\leq s+\log(1+\mathrm{e}^{-s})+\sup_Yv$, so $\Phi_v(s)-s$ is bounded above as $s\to+\infty$, which forces $\Phi_v'\leq1$.

On $\C^k\setminus\{0\}$ the function $v$ is locally bounded, and $\mathrm{MA}^{\mathrm{ext}}_\kappa(v)=(\ddc p)^k$ there. The pole $o$ and the hyperplane $\{z_0=0\}$ are pluripolar, so they are not charged (\cref{thm:arbitrary-extended-ma}). If $\Phi_v$ is smooth, the computation in the proof of \cref{lem:projective-sublogarithmic-density}, with $Q(r)=\Phi_v'(\log r)$, shows that the annulus $\{s_1<s<s_2\}$ has mass $\Phi_v'(s_2)^k-\Phi_v'(s_1)^k$. In general, convolve $\Phi_v$ with an even mollifier $\chi_\varepsilon$. The functions $\Phi_\varepsilon=\Phi_v*\chi_\varepsilon$ are smooth and convex, satisfy $0\leq\Phi_\varepsilon'\leq1$, and decrease to $\Phi_v$ as $\varepsilon\downarrow0$, since $\Phi_v(s+\varepsilon x)+\Phi_v(s-\varepsilon x)$ increases with $\varepsilon$ by convexity. The potentials $\Phi_\varepsilon(\log|z|^2)$ decrease to $p$ and are locally uniformly bounded on $\C^k\setminus\{0\}$, so bounded decreasing continuity (see the remark after \cref{prop:bounded-products}) gives weak convergence of their Monge--Amp\`ere measures to $\sigma=\mathrm{MA}^{\mathrm{ext}}_\kappa(v)$ there. If $s_1<s_2$ are points where $\Phi_v$ is differentiable, then $\Phi_\varepsilon'\to\Phi_v'$ at $s_1$ and $s_2$, and comparing the open annulus $\{s_1<s<s_2\}$ with its compact closure gives
\[
 \sigma(\{s_1<s<s_2\})\leq\Phi_v'(s_2)^k-\Phi_v'(s_1)^k\leq\sigma(\{s_1\leq s\leq s_2\}).
\]
Such points are dense. Letting them decrease to given $s_1'<s_2'$ from above, the open and the closed annuli both converge to $\{s_1'<s\leq s_2'\}$, and right continuity of $\Phi_v'$ gives $\sigma(\{s_1'<s\leq s_2'\})=\Phi_v'(s_2')^k-\Phi_v'(s_1')^k$. Letting $s_1'\to-\infty$ proves \eqref{eq:nonlp-selection-flux}, and letting also $s_2'\to+\infty$ gives the total mass.

Suppose now that $v\geq u_\delta-c$. The function $\Phi_{u_\delta}(s)=\log(1+\mathrm{e}^s)+1-L^\delta$, with $L=1-s+\log(1+\mathrm{e}^s)$, satisfies $\Phi_{u_\delta}(s)/s\to0$ as $s\to-\infty$. By convexity, $\Phi_v(s)\leq\Phi_v(0)+\Phi_v'(-\infty)s$ for $s<0$, while $\Phi_v\geq\Phi_{u_\delta}-c$; hence $\Phi_v'(-\infty)=0$. Near the hyperplane $\{z_0=0\}$ the function $u_\delta$ is bounded, so $v$ is bounded below there, and $\Phi_v(s)-s$ is bounded below as $s\to+\infty$. This function is convex and nonincreasing, so its derivative $\Phi_v'-1$ tends to zero; thus $\Phi_v'(+\infty)=1$, and \eqref{eq:nonlp-selection-flux} gives the formula for $\Phi_v'$. Moreover $\Phi_v(s)-\log(1+\mathrm{e}^s)$ has a finite limit as $s\to+\infty$: it is the sum of the nonincreasing function $\Phi_v(s)-s$, which is bounded below, and of $-\log(1+\mathrm{e}^{-s})$. By \cref{lem:qpsh-removable-analytic,lem:qpsh-almost-everywhere-order}, the value of $v$ at a point of $\{z_0=0\}$ is the upper limit of its values off this hyperplane, which is this limit. Since $\Phi_v$ is continuous, $v$ is continuous on $Y\setminus\{o\}$.

Finally, let $v_1,v_2$ be invariant with $\mathrm{MA}^{\mathrm{ext}}_\kappa(v_1)=\mathrm{MA}^{\mathrm{ext}}_\kappa(v_2)$ of total mass one. Since $\Phi_{v_i}'\leq1$, the total mass forces $\Phi_{v_i}'(-\infty)=0$, and \eqref{eq:nonlp-selection-flux} gives $\Phi_{v_1}'=\Phi_{v_2}'$. Hence $v_1-v_2$ is constant off the origin of the chart, and therefore on $Y$ by \cref{lem:qpsh-almost-everywhere-order}.
\end{proof}

\subsection{Separated greatest solutions}

On the projective factor, with background $\kappa$ and volume form $\kappa^k$, consider the exponential equation
\begin{equation}\label{eq:nonlp-selection-factor-equation}
 \mathrm{MA}^{\mathrm{ext}}_\kappa(w)=\mathrm{e}^{\lambda w}q_\alpha\kappa^k,\qquad w\in\PSH(Y,\kappa),
\end{equation}
for $0<\alpha<1$ and $\lambda>0$. Since $u_\alpha\leq0$ and $\mathrm{MA}^{\mathrm{ext}}_\kappa(u_\alpha)=q_\alpha\kappa^k$, the potential $u_\alpha$ is an exponential subsolution of \eqref{eq:nonlp-selection-factor-equation} for every $\lambda>0$. By \cref{thm:individual-exponential-subsolution}, applied on $Y$, the equation has a greatest solution, which we denote by $w_{\lambda,\alpha}$. Put
\[
 a_{\lambda,\alpha}=\sup_Yw_{\lambda,\alpha},\qquad
 v_{\lambda,\alpha}=w_{\lambda,\alpha}-a_{\lambda,\alpha},
\]
and let $U_{\lambda,\alpha}$ be the greatest solution of $\mathrm{MA}^{\mathrm{ext}}_H(U)=\mathrm{e}^{\lambda U}f_\alpha\,\mathrm{d}V$ on $X$, with $c_{\lambda,\alpha}=\exp(\lambda\sup_XU_{\lambda,\alpha})$ as in \cref{cor:individual-positive-threshold-exponential}.

\Needspace{12\baselineskip}
\begin{proposition}\label{prop:nonlp-selection-separated}\index{Monge--Ampere equation@Monge--Amp\`ere equation!separated greatest solutions}
Let $0<\alpha<1$ and $\lambda>0$. The greatest solution $w_{\lambda,\alpha}$ is invariant and satisfies $w_{\lambda,\alpha}\geq u_\alpha$, and its Monge--Amp\`ere measure is a probability measure. Moreover,
\[
 U_{\lambda,\alpha}=\lambda^{-1}\log b+\pi_Y^*w_{\lambda,\alpha},
\]
the measure $\mathrm{MA}^{\mathrm{ext}}_H(U_{\lambda,\alpha})$ has total mass $b$, and $U_{\lambda,\alpha}-\sup_XU_{\lambda,\alpha}=\pi_Y^*v_{\lambda,\alpha}$. The rescaled scalar error is
\begin{equation}\label{eq:nonlp-selection-scalar-error}
 \lambda^{-1}\log(c_{\lambda,\alpha}/b)=a_{\lambda,\alpha}
 =-\frac1\lambda\log\int_Y\mathrm{e}^{\lambda v_{\lambda,\alpha}}q_\alpha\kappa^k\geq0.
\end{equation}
\end{proposition}
\begin{proof}
Write $w=w_{\lambda,\alpha}$. Since $u_\alpha$ is an exponential subsolution and $w$ is the greatest one, $w\geq u_\alpha$; in particular $a_{\lambda,\alpha}\geq\sup_Yu_\alpha=0$. For $g\in G$, the function $w\circ g$ is again a solution, because $g$ preserves $\kappa$ and $q_\alpha\kappa^k$. Hence $w\circ g\leq w$ for every $g\in G$; applying this to $g^{-1}$ at the points $g\cdot y$ gives the reverse inequality, so $w$ is invariant. By \cref{lem:nonlp-selection-radial}, $\mathrm{MA}^{\mathrm{ext}}_\kappa(w)$ is a probability measure, that is, $\int_Y\mathrm{e}^{\lambda w}q_\alpha\kappa^k=1$.

For every $w'\in\PSH(Y,\kappa)$, the computation in Step~3 of the proof of \cref{thm:closed-factor-product-equation}, which uses only that $w'$ is $\kappa$-plurisubharmonic, gives
\[
 \mathrm{MA}^{\mathrm{ext}}_H(\pi_Y^*w')=Ch_0^m\wedge\pi_Y^*\mathrm{MA}^{\mathrm{ext}}_\kappa(w').
\]
Since $Ch_0^m\wedge\kappa^k=b\mu$ and $q_\alpha\mu=f_\alpha\,\mathrm{d}V$, the function $U'=\lambda^{-1}\log b+\pi_Y^*w$ satisfies
\[
 \mathrm{MA}^{\mathrm{ext}}_H(U')=b\,\mathrm{e}^{\lambda\pi_Y^*w}q_\alpha\mu=\mathrm{e}^{\lambda U'}f_\alpha\,\mathrm{d}V,
\]
so $U'\leq U_{\lambda,\alpha}$ by maximality. Its total mass is $b\int_Y\mathrm{e}^{\lambda w}q_\alpha\kappa^k=b$.

Conversely, by \cref{thm:individual-exponential-subsolution}, $w$ is the decreasing limit of the continuous solutions $w_j$ of $\mathrm{MA}_\kappa(w_j)=\mathrm{e}^{\lambda w_j}\min(q_\alpha,j)\kappa^k$. The continuous functions $W_j=\lambda^{-1}\log b+\pi_Y^*w_j$ satisfy, by the same expansion,
\[
 \mathrm{MA}_H(W_j)=\mathrm{e}^{\lambda W_j}R\min(q_\alpha,j)\,\mathrm{d}V
 \leq\mathrm{e}^{\lambda W_j}f_\alpha\,\mathrm{d}V.
\]
Let $\tau>0$. On the set $\{W_j<U_{\lambda,\alpha}-\tau\}$,
\[
 \mathrm{MA}_H(W_j)\leq\mathrm{e}^{\lambda(U_{\lambda,\alpha}-\tau)}f_\alpha\,\mathrm{d}V
 =\mathrm{e}^{-\lambda\tau}\mathrm{MA}^{\mathrm{ext}}_H(U_{\lambda,\alpha}-\tau).
\]
Since $W_j$ is bounded, $P_{H/2}(W_j)\geq\inf_XW_j$ is not identically $-\infty$, and $U_{\lambda,\alpha}-\tau\preceq W_j$. As in the proof of \cref{thm:individual-exponential-subsolution}, \cref{thm:singular-strict-domination} gives $W_j\geq U_{\lambda,\alpha}-\tau$. Letting $\tau\downarrow0$ and then $j\to\infty$ gives $U_{\lambda,\alpha}\leq U'$, hence $U_{\lambda,\alpha}=U'$.

Consequently $\sup_XU_{\lambda,\alpha}=\lambda^{-1}\log b+a_{\lambda,\alpha}$, which gives $U_{\lambda,\alpha}-\sup_XU_{\lambda,\alpha}=\pi_Y^*v_{\lambda,\alpha}$ and $c_{\lambda,\alpha}=b\,\mathrm{e}^{\lambda a_{\lambda,\alpha}}$. Finally, $1=\int_Y\mathrm{e}^{\lambda w}q_\alpha\kappa^k=\mathrm{e}^{\lambda a_{\lambda,\alpha}}\int_Y\mathrm{e}^{\lambda v_{\lambda,\alpha}}q_\alpha\kappa^k$, which is \eqref{eq:nonlp-selection-scalar-error}.
\end{proof}

By Fubini's theorem, the integral in \eqref{eq:nonlp-selection-scalar-error} equals $\int_X\mathrm{e}^{\lambda(U_{\lambda,\alpha}-\sup_XU_{\lambda,\alpha})}f_\alpha\,\mathrm{d}V$. The identity thus expresses the rescaled scalar error through the damping of the density by the normalized solution, a consequence of the constant mass $b$ of the exact solutions. At zero threshold, by contrast, these masses tend to zero (\cref{prop:zero-threshold-exponential-mass-collapse}).

\subsection{The vanishing exponent}

By \eqref{eq:nonlp-selection-radial-data} and Tonelli's theorem, since $L^\alpha-1=\int_1^L\alpha x^{\alpha-1}\,\mathrm{d}x$,
\begin{equation}\label{eq:nonlp-selection-moment-formula}
 J_\alpha=\alpha\int_1^\infty x^{\alpha-1}Q_\alpha(x)^k\,\mathrm{d}x,\qquad
 Q_\alpha(x)=\alpha x^{\alpha-1}+\bigl(1-\alpha x^{\alpha-1}\bigr)\mathrm{e}^{1-x}.
\end{equation}
As $x\to+\infty$, $Q_\alpha(x)\sim\alpha x^{\alpha-1}$, so the integrand is asymptotic to $\alpha^{k+1}x^{(k+1)(\alpha-1)}$. Hence $J_\alpha<\infty$ if and only if $(k+1)\alpha<k$. Moreover $J_\alpha>0$, since $u_\alpha<0$ off the hyperplane $\{z_0=0\}$.

\Needspace{24\baselineskip}
\begin{theorem}\label{thm:nonlp-selection-limit}\index{Exponential approximation!product densities}\index{Normalization scalar!rescaled scalar error}
Let $0<\alpha<1$, and let $\alpha_j\in(0,1)$ and $\lambda_j>0$ satisfy $\alpha_j\to\alpha$ and $\lambda_j\to0$, with no condition relating the two rates. Write
\[
 U_j=U_{\lambda_j,\alpha_j},\quad w_j=w_{\lambda_j,\alpha_j},\quad a_j=a_{\lambda_j,\alpha_j},\quad v_j=v_{\lambda_j,\alpha_j},\quad c_j=c_{\lambda_j,\alpha_j}.
\]
\begin{enumerate}
\item The densities $q_{\alpha_j}$ converge to $q_\alpha$ in $L^1(\kappa^k)$, the scalars satisfy $c_j\to b$, and $\mathrm{MA}^{\mathrm{ext}}_\kappa(v_j)\to q_\alpha\kappa^k$ in total variation.
\item The normalized potentials $v_j$ converge to $u_\alpha$ in $L^1(Y)$, in capacity for $\operatorname{Cap}_\kappa$, and uniformly on compact subsets of $Y\setminus\{o\}$. Consequently $U_j-\sup_XU_j\to\pi_Y^*u_\alpha$ in $L^1(X)$, in capacity, and uniformly on compact subsets of $X\setminus(X_0\times\{o\})$.
\item If $(k+1)\alpha<k$, then the rescaled scalar errors $\lambda_j^{-1}\log(c_j/b)=a_j$ converge to $J_\alpha$, and
\[
 U_j-\lambda_j^{-1}\log b=\pi_Y^*w_j\longrightarrow\pi_Y^*(u_\alpha+J_\alpha)
\]
in the three senses of \textup{(2)}. The limit satisfies $\mathrm{MA}^{\mathrm{ext}}_H(\pi_Y^*(u_\alpha+J_\alpha))=bf_\alpha\,\mathrm{d}V$ and $\int_X\pi_Y^*(u_\alpha+J_\alpha)f_\alpha\,\mathrm{d}V=0$.
\item If $(k+1)\alpha\geq k$, then $a_j\to+\infty$ while $\lambda_ja_j\to0$, and $U_j-\lambda_j^{-1}\log b\to+\infty$ uniformly on compact subsets of $X\setminus(X_0\times\{o\})$. This holds also when $(k+1)\alpha_j<k$ for every $j$.
\end{enumerate}
\end{theorem}
\begin{proof}
Every subsequence has a further subsequence along which $\lambda_j$ decreases, and all assertions are statements of convergence; we may therefore assume that $\lambda_j\downarrow0$. Choose $\alpha_+$ with $\alpha<\alpha_+<1$, and with $(k+1)\alpha_+<k$ if $(k+1)\alpha<k$. Discarding finitely many indices, we may assume $\alpha_j\leq\alpha_+$.

\emph{Step 1: a common dominating density.} For $L\geq1$ and $0<\delta\leq\alpha_+$ we have $\delta L^{\delta-1}\leq L^{\alpha_+-1}$, $\delta(1-\delta)L^{\delta-2}\leq L^{\alpha_+-2}$ and $0<1-\delta L^{\delta-1}\leq1$, while $\mathrm{e}^{1-L}L^2$ is bounded. Hence $Q_\delta\leq C_1L^{\alpha_+-1}$ and $D_\delta=t(1-\delta L^{\delta-1})+(1-t)\delta(1-\delta)L^{\delta-2}\leq C_1L^{\alpha_+-2}$, with $C_1$ independent of $\delta$. By \eqref{eq:nonlp-selection-radial-data}, every $q_\delta$ with $0<\delta\leq\alpha_+$ is bounded by
\[
 g=C_1^kL^{-1-k(1-\alpha_+)}\mathrm{e}^{k(L-1)}.
\]
The function $g$ is bounded on every $Y_T$, and since $\kappa^k$ has radial distribution $k\mathrm{e}^{-k(L-1)}\,\mathrm{d}L$,
\begin{gather*}
 \int_Yg\kappa^k=kC_1^k\int_1^\infty L^{-1-k(1-\alpha_+)}\,\mathrm{d}L<\infty,\\
 \int_Y(-u_{\alpha_+})g\kappa^k=kC_1^k\int_1^\infty(L^{\alpha_+}-1)L^{-1-k(1-\alpha_+)}\,\mathrm{d}L,
\end{gather*}
and the second integral is finite when $(k+1)\alpha_+<k$. The functions $q_\delta$ depend continuously on $\delta$ at every point of $Y\setminus\{o\}$, so dominated convergence gives $q_{\alpha_j}\to q_\alpha$ in $L^1(\kappa^k)$.

\emph{Step 2: scalars and measures.} The measure $g\kappa^k$ is finite and does not charge pluripolar sets, so \cref{lem:normalized-exponential-npp}, applied on $Y$, gives
\[
 \varepsilon_j=\int_Y\bigl(1-\mathrm{e}^{\lambda_jv_j}\bigr)g\kappa^k\longrightarrow0.
\]
By \eqref{eq:nonlp-selection-scalar-error}, $\mathrm{e}^{-\lambda_ja_j}=\int_Y\mathrm{e}^{\lambda_jv_j}q_{\alpha_j}\kappa^k\geq1-\varepsilon_j$. Since $a_j\geq0$, this gives $\lambda_ja_j\to0$, hence $c_j=b\,\mathrm{e}^{\lambda_ja_j}\to b$. The measures $\mathrm{MA}^{\mathrm{ext}}_\kappa(v_j)=\mathrm{e}^{\lambda_ja_j}\mathrm{e}^{\lambda_jv_j}q_{\alpha_j}\kappa^k$ are probability measures, and
\[
 \bigl\|\mathrm{MA}^{\mathrm{ext}}_\kappa(v_j)-q_\alpha\kappa^k\bigr\|_{\mathrm{TV}}
 \leq\bigl(\mathrm{e}^{\lambda_ja_j}-1\bigr)+\varepsilon_j+\|q_{\alpha_j}-q_\alpha\|_{L^1(\kappa^k)}
 \longrightarrow0.
\]
This proves (1).

\emph{Step 3: radial derivatives.} By \cref{prop:nonlp-selection-separated}, $v_j$ is invariant and $v_j\geq u_{\alpha_j}-a_j$. With the notation of \cref{lem:nonlp-selection-radial}, put $\Phi_j=\Phi_{v_j}$ and $\Phi=\Phi_{u_\alpha}$, and let $F_j(s)$ and $F(s)$ be the masses of $\{|z|^2\leq\mathrm{e}^s\}$ for $\mathrm{MA}^{\mathrm{ext}}_\kappa(v_j)$ and for $q_\alpha\kappa^k$. The lemma gives $\Phi_j'=F_j^{1/k}$ and $\Phi'=F^{1/k}$, and Step~2 gives $\sup_s|F_j(s)-F(s)|\leq\eta_j\to0$, where $\eta_j$ is the total-variation distance in (1). Since $|x^{1/k}-y^{1/k}|\leq|x-y|^{1/k}$ for $x,y\geq0$,
\[
 \sup_{s\in\mathbb R}|\Phi_j'(s)-\Phi'(s)|\leq\eta_j^{1/k}\longrightarrow0.
\]
For $s\geq0$, the set $\{|z|^2>\mathrm{e}^s\}$ lies in $Y_2$, where $g\leq C_1^k\mathrm{e}^k$, and its $\kappa^k$-measure is $1-(1+\mathrm{e}^{-s})^{-k}\leq k\mathrm{e}^{-s}$. Since $v_j\leq0$ and $\mathrm{e}^{\lambda_ja_j}\leq2$ for large $j$, we obtain $1-F_j(s)\leq C_2\mathrm{e}^{-s}$, and likewise $1-F(s)\leq C_2\mathrm{e}^{-s}$. For $s\geq s_0=\log(2C_2)$ both $F_j$ and $F$ lie in $[1/2,1]$, where $x\mapsto x^{1/k}$ is $2$-Lipschitz, so $|\Phi_j'(s)-\Phi'(s)|\leq4C_2\mathrm{e}^{-s}$. Splitting at a large $S$, we conclude that
\[
 \int_A^\infty|\Phi_j'(s)-\Phi'(s)|\,\mathrm{d}s\leq(S-A)\eta_j^{1/k}+4C_2\mathrm{e}^{-S}
\]
for $S\geq\max(A,s_0)$, so the left side tends to zero for every $A\in\mathbb R$.

\emph{Step 4: the normalized potentials.} By \cref{lem:nonlp-selection-radial}, $v_j$ and $u_\alpha$ are continuous on $Y\setminus\{o\}$, and on the chart $v_j-u_\alpha=\Phi_j(s)-\Phi(s)$. For $T>1$, $Y_T$ is the closure in $Y$ of a set $\{s\geq A\}$, so Step~3 shows that the oscillation of $v_j-u_\alpha$ on $Y_T$ tends to zero; since the sets $Y_T$ increase with $T$, this holds for every $T\geq1$. Fix a point $y_0$ of the hyperplane $\{z_0=0\}$ and put $d_j=(v_j-u_\alpha)(y_0)$; thus $v_j-u_\alpha-d_j\to0$ uniformly on every $Y_T$. Given a subsequence, \cref{thm:normalized-compactness} gives a further subsequence along which $v_j\to v$ in $L^1(Y)$, with $\sup_Yv=0$. Along it, integrating over $Y_T$ shows that $d_j$ converges to some $d$ and that $v=u_\alpha+d$ almost everywhere on every $Y_T$, hence on $Y$ by \cref{lem:qpsh-almost-everywhere-order}. Then $d=\sup_Yv-\sup_Yu_\alpha=0$. Consequently $v_j\to u_\alpha$ in $L^1(Y)$ and uniformly on every $Y_T$, along the whole sequence.

For convergence in capacity, let $\varepsilon>0$, and let $T$ be so large that $T^\alpha>2$. For large $j$,
\[
 \{|v_j-u_\alpha|>\varepsilon\}\subset Y\setminus Y_T=\{u_\alpha<1-T^\alpha\},
\]
and \cref{lem:global-capacity-estimates}, applied on $Y$ to the normalized potential $u_\alpha$, bounds the capacity of the right side by a constant multiple of $(T^\alpha-1)^{-1}$. On $X$, $U_j-\sup_XU_j-\pi_Y^*u_\alpha=\pi_Y^*(v_j-u_\alpha)$ by \cref{prop:nonlp-selection-separated}. Since $\mathrm{d}V\leq(\min_XR)^{-1}\mu$ and $\int_X|\pi_Y^*\xi|\,\mu=\int_Y|\xi|\kappa^k$ for integrable $\xi$ on $Y$, $L^1$ convergence passes to $X$; every compact subset of $X\setminus(X_0\times\{o\})$ lies in some $\pi_Y^{-1}(Y_T)$; and the exceptional sets above pull back to sublevel sets of the normalized $H$-plurisubharmonic function $\pi_Y^*u_\alpha$. If $\omega$ is the Hermitian form defining the capacity on $X$, then after division by a constant $K$ with $H\leq K\omega$, the same lemma bounds their capacity by a constant multiple of $(T^\alpha-1)^{-1}$. This proves (2).

\emph{Step 5: finite first moment.} Suppose $(k+1)\alpha<k$, so that $(k+1)\alpha_+<k$ and $C_3=\int_Y(-u_{\alpha_+})g\kappa^k$ is finite. Write $\nu_j=q_{\alpha_j}\kappa^k$, a probability measure. Jensen's inequality and \eqref{eq:nonlp-selection-scalar-error} give $\mathrm{e}^{-\lambda_ja_j}\geq\exp(\lambda_j\int_Yv_j\,\mathrm{d}\nu_j)$, that is, $a_j\leq\int_Y(-v_j)\,\mathrm{d}\nu_j$. Choose $T\geq2$ with $\int_{Y\setminus Y_T}g\kappa^k\leq1/2$. On $Y_T$, Step~4 gives $v_j\geq u_\alpha-1\geq-T^\alpha$ for large $j$. Off $Y_T$, we use $v_j=w_j-a_j\geq u_{\alpha_j}-a_j\geq u_{\alpha_+}-a_j$, since $L^{\alpha_j}\leq L^{\alpha_+}$. Hence
\[
 a_j\leq T^\alpha+C_3+a_j\nu_j(Y\setminus Y_T)\leq T^\alpha+C_3+a_j/2,
\]
so $0\leq a_j\leq C_4=2(T^\alpha+C_3)$. Now put
\[
 I_j=\int_Y\frac{1-\mathrm{e}^{\lambda_jv_j}}{\lambda_j}\,\mathrm{d}\nu_j,
 \qquad\text{so that}\qquad
 \mathrm{e}^{-\lambda_ja_j}=1-\lambda_jI_j.
\]
The integrands $\lambda_j^{-1}(1-\mathrm{e}^{\lambda_jv_j})q_{\alpha_j}$ lie between $0$ and $(-v_j)q_{\alpha_j}\leq(-u_{\alpha_+}+C_4)g$, which is integrable, and they converge to $(-u_\alpha)q_\alpha$ on $Y\setminus\{o\}$ by Step~4 and the pointwise convergence of $q_{\alpha_j}$. Dominated convergence gives $I_j\to J_\alpha$. Since $x\leq-\log(1-x)\leq x/(1-x)$ for $0\leq x<1$,
\[
 I_j\leq a_j=-\lambda_j^{-1}\log(1-\lambda_jI_j)\leq\frac{I_j}{1-\lambda_jI_j},
\]
and $\lambda_jI_j=1-\mathrm{e}^{-\lambda_ja_j}\to0$; hence $a_j\to J_\alpha$. By \cref{prop:nonlp-selection-separated}, $U_j-\lambda_j^{-1}\log b=\pi_Y^*w_j=\pi_Y^*(v_j+a_j)$, and (2) gives the asserted convergence. The equation for the limit is that of \cref{thm:zero-lower-mass-product-equation}, since adding a constant does not change the measure, and the integral vanishes by \eqref{eq:nonlp-selection-moment}.

\emph{Step 6: infinite first moment.} Suppose $(k+1)\alpha\geq k$, so that $J_\alpha=+\infty$. Define $I_j$ as in Step~5; the identity $\mathrm{e}^{-\lambda_ja_j}=1-\lambda_jI_j$ and the inequality $a_j\geq I_j$ still hold. The integrands are nonnegative and converge to $(-u_\alpha)q_\alpha$ on $Y\setminus\{o\}$, so Fatou's lemma gives $\liminf_jI_j\geq J_\alpha=+\infty$. Hence $a_j\to+\infty$, while $\lambda_ja_j\to0$ by Step~2. Since $v_j\to u_\alpha$ uniformly on every $Y_T$, where $u_\alpha$ is bounded, $w_j=v_j+a_j\to+\infty$ uniformly there. This argument does not use the position of $\alpha_j$ relative to $k/(k+1)$.
\end{proof}

Taking $\alpha_j=\alpha$ for all $j$, the theorem applies to every sequence $\lambda_j\to0$, and hence to the whole family: as $\lambda\downarrow0$, the normalized greatest solutions $U_{\lambda,\alpha}-\sup_XU_{\lambda,\alpha}$ converge to $\pi_Y^*u_\alpha$, and the rescaled scalar error $\lambda^{-1}\log(c_{\lambda,\alpha}/b_H(f_\alpha))$ converges to the explicit constant \eqref{eq:nonlp-selection-moment-formula} if $(k+1)\alpha<k$, and tends to $+\infty$ otherwise. In the second case the scalars still converge, $c_{\lambda,\alpha}\to b_H(f_\alpha)$, but not at the rate $\mathrm{e}^{O(\lambda)}$.

The limit in part (3) is the potential selected by requiring that its mean against $f_\alpha\,\mathrm{d}V$ vanish, and the normalizing probability is $f_\alpha\,\mathrm{d}V$ itself. It also has an order description among invariant potentials. By \cref{thm:closed-factor-product-equation}, applied with the density $q_\alpha$ on the projective factor, every $v\in\PSH(Y,\kappa)$ with $\mathrm{MA}^{\mathrm{ext}}_\kappa(v)\geq q_\alpha\kappa^k$ satisfies equality and $\mathrm{MA}^{\mathrm{ext}}_H(\pi_Y^*v)=bf_\alpha\,\mathrm{d}V$. If $v$ is invariant, \cref{lem:nonlp-selection-radial} gives $v=u_\alpha+d$ for a constant $d$. Conversely, for a measure $\sigma$ on $Y$, evaluation on the sets $X_0\times E$ shows that $Ch_0^m\wedge\pi_Y^*\sigma\geq bf_\alpha\,\mathrm{d}V$ exactly when $\sigma\geq q_\alpha\kappa^k$. The threshold subsolutions pulled back from invariant functions on $Y$ are therefore exactly the functions $\pi_Y^*(u_\alpha+d)$, and in the first-moment regime $\pi_Y^*(u_\alpha+J_\alpha)$ is the greatest of them whose mean against $f_\alpha\,\mathrm{d}V$ is nonpositive. This is the form of the selection in \cref{thm:vanishing-exponential-selection} for $L^p$ data. When $(k+1)\alpha\geq k$, every translate of $\pi_Y^*u_\alpha$ has $f_\alpha$-mean $-\infty$, so this normalization selects nothing, and the shifted solutions escape to $+\infty$. On $Y$ itself, $u_\alpha$ is the greatest normalized threshold subsolution of \cref{prop:greatest-normalized-threshold} for the density $q_\alpha$: that function is invariant, because $G$ preserves the family of normalized threshold subsolutions, so by the preceding argument it is a translate of $u_\alpha$, and both have supremum zero.

The divergence in part (4) also bounds what uniform root tails can give. By \cref{prop:closed-factor-root-subspace}, the root $p_\alpha=f_\alpha^{1/N}$ belongs to $\mathcal D_H^0$ for every $0<\alpha<1$. If the modulus $\eta_{p_\alpha}(\lambda)$ of \cref{prop:exponential-root-tail-modulus} were $O(\lambda)$, the estimate following that proposition would keep the rescaled scalar error bounded. For $(k+1)\alpha\geq k$ it tends to $+\infty$; hence $\eta_{p_\alpha}(\lambda)/\lambda$ is unbounded as $\lambda\downarrow0$, and uniform root tails do not control the rescaled scalar error.

Part (4) includes parameters $\alpha_j<k/(k+1)$ that increase to the moment boundary. For each fixed $j$ the rescaled scalar errors converge to $J_{\alpha_j}$ as $\lambda\downarrow0$, and $J_{\alpha_j}\to+\infty$ by Fatou's lemma; the theorem shows that the joint limit diverges for every relation between the two rates. Throughout, the reference is zero, and no singular reference is varied.

\Needspace{10\baselineskip}
\begin{corollary}\label{cor:nonlp-selection-fivefold}\index{Elliptic fibration!zero-exponent limit}
On the fivefold $X_0\times\mathbb P^2$ of \cref{cor:elliptic-fivefold-exact-threshold}, the normalized greatest solutions of $\mathrm{MA}^{\mathrm{ext}}_H(U)=\mathrm{e}^{\lambda U}f_\alpha\,\mathrm{d}V$ converge to $\pi_Y^*u_\alpha$ in $L^1$, in capacity and uniformly on compact subsets of the complement of the pole fiber as $\lambda\downarrow0$, and $c_{\lambda,\alpha}\to240\pi$. The rescaled scalar error $\lambda^{-1}\log(c_{\lambda,\alpha}/240\pi)$ converges to $J_\alpha$ if $\alpha<2/3$, and tends to $+\infty$ if $\alpha\geq2/3$.
\end{corollary}
\begin{proof}
Here $m=3$, $k=2$ and $b=240\pi$ by \cref{cor:elliptic-fivefold-exact-threshold}, and $(k+1)\alpha<k$ means $\alpha<2/3$. Apply \cref{thm:nonlp-selection-limit} with $\alpha_j=\alpha$ along every sequence $\lambda_j\to0$.
\end{proof}

These conclusions rest on the product structure, on the invariance under $G$ and on the radial form of the equation on the projective factor. For a general integrable density with positive threshold at zero lower mass, the convergence of the whole family of normalized greatest solutions and the behavior of the rescaled scalar error remain open. Even for the densities $f_\alpha$, we do not know whether $\pi_Y^*u_\alpha$ is the greatest normalized threshold subsolution on $X$, nor whether the selected limit has an order description among threshold subsolutions that are not pulled back from the projective factor.
